\apisummary{
  Collectively allocate a zeroed block of symmetric memory in a specific space.
}

\begin{apidefinition}

\begin{Csynopsis}
void *@\FuncDecl{shmem\_space\_calloc}@(shmem_space_t space, size_t count, size_t size);
\end{Csynopsis}

\begin{apiarguments}
  \apiargument{IN}{space}{The space within which to allocate.}
  \apiargument{IN}{count}{The number of elements to allocate.}
  \apiargument{IN}{size}{The size in bytes of each element to allocate.}
\end{apiarguments}

\apidescription{
  The \FUNC{shmem\_space\_calloc} routine is a collective operation on the
  team associated with \VAR{space}. It allocates a region of
  remotely-accessible symmetric memory for an array of \VAR{count} objects of
  \VAR{size} bytes each from the symmetric heap identified by \VAR{space}, and
  returns the symmetric address of the lowest byte of the allocated symmetric
  memory. The space is initialized to all bits zero.

  If the allocation succeeds, the pointer returned shall be suitably aligned so
  that it may be assigned to a pointer to any type of object. If the allocation
  does not succeed, or either \VAR{count} or \VAR{size} is \CONST{0}, the
  return value is a null pointer.

  The values of the \VAR{space}, \VAR{count}, and \VAR{size} arguments must be
  identical on all \acp{PE}; otherwise, the behavior is undefined.

  When \VAR{count} or \VAR{size} is \CONST{0}, the \FUNC{shmem\_space\_calloc}
  routine returns without performing a barrier; otherwise,
  \FUNC{shmem\_space\_calloc} calls a procedure that is semantically
  equivalent to \FUNC{shmem\_team\_sync} on exit.
}

\apireturnvalues{
  The \FUNC{shmem\_space\_calloc} routine returns the symmetric address of the
  allocated space; otherwise, it returns a null pointer.
}

\end{apidefinition}


